% GENERATED FILE. Edit canonical lessons, then run npm run generate:books.
% canonical-source-hash: fnv1a64:e92969633a0909d7
% canonical-lessons: AR-C350-nashir, AR-C350-taba, AR-C350-nuskha, AR-C350-manshur, AR-C350-mulhaq

\chapter{Publishing}
\label{ch:ar-publishing}
\hlchaptermodality{350}

\begin{chapteropening}
\textbf{By the end of this chapter:} \emph{I can say a publisher, an edition, a copy, a leaflet and a supplement.}

Complete the last lesson of chapter 350: I can say a publisher, an edition, a copy, a leaflet and a supplement.
\end{chapteropening}

\section[\texorpdfstring{\ar{ناشر} (nāshir)}{ناشر (nāshir)}]{\ar{ناشر} (nāshir) — a publisher}

\label{lesson:AR-C350-nashir}



\begin{quote}
Before the new one: say the Arabic for someday, then the Arabic for right after.
\end{quote}

\subsection*{You'll want to know: \ar{ناشر}}
\textbf{\ar{ناشر}} — \emph{nāshir} — \textquotedblleft{}a publisher\textquotedblright{}.

\begin{quote}
\ar{قبل} \ar{الناشر} \ar{كتاب} \ar{أختي}.
\end{quote}

\begin{quote}
\emph{qabila an-nāshir kitāb ukhtī.}
\end{quote}

\begin{quote}
The publisher accepted my sister’s book.
\end{quote}

A publishing house is \ar{دار} \ar{نشر}.

\begin{cousinweb}
\textbf{\ar{ناشر}} (\emph{nāshir}) is built on the native Arabic root \emph{n-sh-r} (Lane), \textquotedblleft{}to spread out\textquotedblright{}, the root of the taught \emph{nashara}. \emph{Nāshir} is the active participle, \textquotedblleft{}one who publishes\textquotedblright{}.
\end{cousinweb}

\begin{grammarlens}[title={What you've built}]
One more word for telling what happened.
\end{grammarlens}

\subsection*{Your turn}
\begin{itemize}
\raggedright
  \item \emph{Say it:} \emph{nāshir}
  \item \emph{Say it:} \emph{nāshir}, once more
  \item \emph{Say it:} \emph{qabila an-nāshir kitāb ukhtī.} — The publisher accepted my sister’s book.
  \item \emph{Recall:} say the Arabic for someday, then the Arabic for right after, then say \emph{nāshir} again
\end{itemize}

\begin{tcolorbox}[breakable,colback=teal!4,colframe=teal!35!black,title={Before you move on}]
What does \ar{ناشر} mean? (\textquotedblleft{}A publisher\textquotedblright{}.) Say it once more.
\end{tcolorbox}
\section[\texorpdfstring{\ar{طبعة} (ṭabʿa)}{طبعة (ṭabʿa)}]{\ar{طبعة} (ṭabʿa) — an edition, a printing}

\label{lesson:AR-C350-taba}



\begin{quote}
Before the new one: say the Arabic for right after, then the Arabic for a publisher.
\end{quote}

\subsection*{You'll want to know: \ar{طبعة}}
\textbf{\ar{طبعة}} — \emph{ṭabʿa} — \textquotedblleft{}an edition, a printing\textquotedblright{}.

\begin{quote}
\ar{صدرت} \ar{طبعة} \ar{جديدة} \ar{من} \ar{الكتاب}.
\end{quote}

\begin{quote}
\emph{ṣadarat ṭabʿa jadīda min al-kitāb.}
\end{quote}

\begin{quote}
A new edition of the book came out.
\end{quote}

\begin{cousinweb}
\textbf{\ar{طبعة}} (\emph{ṭabʿa}) is built on the native Arabic root \emph{ṭ-b-ʿ} (Lane), \textquotedblleft{}to seal, to stamp\textquotedblright{}, the root of the taught \emph{ṭabaʿa}, to print. \emph{Ṭabʿa} is one printing.
\end{cousinweb}

\begin{grammarlens}[title={What you've built}]
One more word for telling what happened.
\end{grammarlens}

\subsection*{Your turn}
\begin{itemize}
\raggedright
  \item \emph{Say it:} \emph{ṭabʿa}
  \item \emph{Say it:} \emph{ṭabʿa}, once more
  \item \emph{Say it:} \emph{ṣadarat ṭabʿa jadīda min al-kitāb.} — A new edition of the book came out.
  \item \emph{Recall:} say the Arabic for right after, then the Arabic for a publisher, then say \emph{ṭabʿa} again
\end{itemize}

\begin{tcolorbox}[breakable,colback=teal!4,colframe=teal!35!black,title={Before you move on}]
What does \ar{طبعة} mean? (\textquotedblleft{}An edition, a printing\textquotedblright{}.) Say it once more.
\end{tcolorbox}
\section[\texorpdfstring{\ar{نسخة} (nuskha)}{نسخة (nuskha)}]{\ar{نسخة} (nuskha) — a printed copy}

\label{lesson:AR-C350-nuskha}



\begin{quote}
Before the new one: say the Arabic for a publisher, then the Arabic for an edition.
\end{quote}

\subsection*{You'll want to know: \ar{نسخة}}
\textbf{\ar{نسخة}} — \emph{nuskha} — \textquotedblleft{}a printed copy\textquotedblright{}.

\begin{quote}
\ar{اشتريت} \ar{نسخة} \ar{من} \ar{الصحيفة}.
\end{quote}

\begin{quote}
\emph{ishtaraytu nuskha min aṣ-ṣaḥīfa.}
\end{quote}

\begin{quote}
I bought a copy of the newspaper.
\end{quote}

Plural: \ar{نسخ}, \emph{nusakh}.

\begin{cousinweb}
\textbf{\ar{نسخة}} (\emph{nuskha}) is built on the native Arabic root \emph{n-s-kh} (Lane), \textquotedblleft{}to copy out a writing\textquotedblright{}, the root of the taught \emph{nasakha}. \emph{Nuskha} is one copy.
\end{cousinweb}

\begin{grammarlens}[title={What you've built}]
One more word for telling what happened.
\end{grammarlens}

\subsection*{Your turn}
\begin{itemize}
\raggedright
  \item \emph{Say it:} \emph{nuskha}
  \item \emph{Say it:} \emph{nuskha}, once more
  \item \emph{Say it:} \emph{ishtaraytu nuskha min aṣ-ṣaḥīfa.} — I bought a copy of the newspaper.
  \item \emph{Recall:} say the Arabic for a publisher, then the Arabic for an edition, then say \emph{nuskha} again
\end{itemize}

\begin{tcolorbox}[breakable,colback=teal!4,colframe=teal!35!black,title={Before you move on}]
What does \ar{نسخة} mean? (\textquotedblleft{}A printed copy\textquotedblright{}.) Say it once more.
\end{tcolorbox}
\section[\texorpdfstring{\ar{منشور} (manshūr)}{منشور (manshūr)}]{\ar{منشور} (manshūr) — a leaflet; a post online}

\label{lesson:AR-C350-manshur}



\begin{quote}
Before the new one: say the Arabic for an edition, then the Arabic for a printed copy.
\end{quote}

\subsection*{You'll want to know: \ar{منشور}}
\textbf{\ar{منشور}} — \emph{manshūr} — \textquotedblleft{}a leaflet; a post online\textquotedblright{}.

\begin{quote}
\ar{وزع} \ar{الطلاب} \ar{منشورا} \ar{عن} \ar{الحفلة}.
\end{quote}

\begin{quote}
\emph{wazzaʿa aṭ-ṭullāb manshūran ʿan al-ḥafla.}
\end{quote}

\begin{quote}
The students handed out a leaflet about the party.
\end{quote}

On social media a \ar{منشور} is a post.

\begin{cousinweb}
\textbf{\ar{منشور}} (\emph{manshūr}) is built on the native Arabic root \emph{n-sh-r} (Lane), \textquotedblleft{}to spread out\textquotedblright{}. \emph{Manshūr} is the passive participle, \textquotedblleft{}a thing spread about\textquotedblright{}.
\end{cousinweb}

\begin{grammarlens}[title={What you've built}]
One more word for telling what happened.
\end{grammarlens}

\subsection*{Your turn}
\begin{itemize}
\raggedright
  \item \emph{Say it:} \emph{manshūr}
  \item \emph{Say it:} \emph{manshūr}, once more
  \item \emph{Say it:} \emph{wazzaʿa aṭ-ṭullāb manshūran ʿan al-ḥafla.} — The students handed out a leaflet about the party.
  \item \emph{Recall:} say the Arabic for an edition, then the Arabic for a printed copy, then say \emph{manshūr} again
\end{itemize}

\begin{tcolorbox}[breakable,colback=teal!4,colframe=teal!35!black,title={Before you move on}]
What does \ar{منشور} mean? (\textquotedblleft{}A leaflet; a post online\textquotedblright{}.) Say it once more.
\end{tcolorbox}
\section[\texorpdfstring{\ar{ملحق} (mulḥaq)}{ملحق (mulḥaq)}]{\ar{ملحق} (mulḥaq) — a supplement; an annex}

\label{lesson:AR-C350-mulhaq}



\begin{quote}
Before the new one: say the Arabic for a printed copy, then the Arabic for a leaflet.
\end{quote}

\subsection*{You'll want to know: \ar{ملحق}}
\textbf{\ar{ملحق}} — \emph{mulḥaq} — \textquotedblleft{}a supplement; an annex\textquotedblright{}.

\begin{quote}
\ar{نشرت} \ar{الصحيفة} \ar{ملحقا} \ar{عن} \ar{الرياضة}.
\end{quote}

\begin{quote}
\emph{nasharat aṣ-ṣaḥīfa mulḥaqan ʿan ar-riyāḍa.}
\end{quote}

\begin{quote}
The newspaper published a sports supplement.
\end{quote}

At an embassy, a \ar{ملحق} is an attaché.

\begin{cousinweb}
\textbf{\ar{ملحق}} (\emph{mulḥaq}) is built on the native Arabic root \emph{l-ḥ-q} (Lane), \textquotedblleft{}to overtake, to join\textquotedblright{}. \emph{Mulḥaq} is the Form IV passive participle, \textquotedblleft{}a thing joined on\textquotedblright{}.
\end{cousinweb}

\begin{grammarlens}[title={What you've built}]
One more word for telling what happened.
\end{grammarlens}

\subsection*{Your turn}
\begin{itemize}
\raggedright
  \item \emph{Say it:} \emph{mulḥaq}
  \item \emph{Say it:} \emph{mulḥaq}, once more
  \item \emph{Say it:} \emph{nasharat aṣ-ṣaḥīfa mulḥaqan ʿan ar-riyāḍa.} — The newspaper published a sports supplement.
  \item \emph{Recall:} say the Arabic for a printed copy, then the Arabic for a leaflet, then say \emph{mulḥaq} again
\end{itemize}

\begin{tcolorbox}[breakable,colback=teal!4,colframe=teal!35!black,title={Before you move on}]
What does \ar{ملحق} mean? (\textquotedblleft{}A supplement; an annex\textquotedblright{}.) Say it once more.
\end{tcolorbox}
